\chapter{Методологија}
\label{ch:metodologija}

Ово поглавље описује поступак којим је спроведен експеримент. Најпре је дат преглед
целокупног истраживачког тока, затим извор и припрема података, потом дизајн
експеримента са свим фиксираним параметрима и, на крају, мере, правила искључивања и
план статистичке анализе. Све вредности наведене у овом поглављу идентичне су у свим
условима експеримента, осим оних које саме представљају експерименталне факторе.

\section{Преглед истраживачког поступка}

Истраживање је спроведено као факторски експеримент у коме се укрштају две
експерименталне осе описане у поглављу~\ref{ch:teorija}: дизајн мерења и архитектура
модела. Циљ поступка јесте да се разлике у тачности предвиђања припишу појединачним
изворима, па су сви кораци подређени захтеву да свака комбинација буде обрађена под
идентичним условима.

Поступак се одвија у четири узастопне фазе:

\begin{enumerate}
\item \textbf{Прикупљање и чишћење података.} Интрадневне свећице на једноминутној
      резолуцији преузимају се из берзанског извора, филтрирају према стварном
      трговачком календару и проверавају на потпуност.
\item \textbf{Конструкција панела мерења.} За сваку од тридесет шест комбинација
      дизајна мерења рачунају се реализоване мере, спроводи се условна декомпозиција
      на континуирану и скоковиту компоненту и граде се HAR агрегати.
\item \textbf{Моделовање.} Сваки модел се оцењује и предвиђа кроз усидрену унапредну
      валидацију, за сваки хоризонт и сваки дизајн од кога зависи.
\item \textbf{Евалуација.} Прогнозе се оцењују функцијама губитка, пореде тестовима
      значајности, подвргавају економској провери и разлажу декомпозицијом варијансе.
\end{enumerate}

Између друге и треће фазе стоји \textbf{замрзнут уговор о шеми панела}: скуп колона,
њихових типова и јединица фиксиран је унапред и третира се као граница интерфејса.
Тиме је обезбеђено да сваки модел види улазне податке у идентичном облику, без обзира
на то из ког дизајна мерења они потичу, што је предуслов за поштено поређење.

\section{Подаци}

\subsection{Извор података}

Полазна напомена тиче се доступности извора. Збирка \textit{Oxford-Man Realized
Library}, која је дуго била стандардни извор у овој области, више није доступна и
институција је саопштила да не планира њену замену. Пошто готово сви радови на које
се овај рад позива користе управо њу, њихова поставка се не може непосредно
поновити, па је извор података морао бити изабран самостално.

Интрадневне свећице на једноминутној резолуцији преузимају се масовним извозом из
берзанског извора. Преузимање се обавља у деловима од по осам година, јер извоз тихо
скраћује одговор на приближно 2,5 милиона редова, што би при шеснаестогодишњем
периоду довело до нечујног губитка података. Свака преузета датотека прати се
\textbf{манифестом порекла} који садржи контролну суму садржаја, број редова и време
преузимања, чиме се обезбеђује поновљивост: поновно преузимање не мења ништа осим ако
се то изричито затражи.

Индекс имплицитне волатилности преузима се из јавне базе економских временских серија
и користи се искључиво са доцњом од једног дана.

\subsection{Узорак}

Узорак обухвата \textbf{тридесет акција} којима се тргује на америчким берзама, у
периоду од 1. јануара 2010. до 31. децембра 2025. године. Симболи су изабрани према
дужини расположиве историје и обиму трговања, уз захтев да свака серија има најмање
шеснаест година једноминутних података.

Панел уместо једне серије изабран је из три разлога. Прво, једна серија од шеснаест
година даје око четири хиљаде дневних опсервација, што је недовољно за стабилно
оцењивање трансформерских архитектура, које су у литератури тестиране на скуповима са
стотинама хиљада тачака, док панел даје ред величине више. Друго, налаз потврђен на
тридесет симбола робуснији је од налаза на једној серији. Треће, панел омогућава
проверу да ли закључци зависе од ликвидности и сектора.

Пет од тридесет симбола односи се на компаније које нису основане у Сједињеним
Државама, али се њиховим акцијама тргује на америчким берзама. То је наведено
ограничење, а не замена за неамеричко тржиште: сви симболи деле исти трговачки
календар, па рад не проверава робусност налаза на тржишту са другачијим радним
временом.

\subsection{Чишћење и контрола квалитета}

Свакој свећици приступа се кроз низ провера изведених из \textbf{стварног трговачког
календара}, а не из фиксираног временског прозора. Фиксиран прозор би на данима са
ранијим затварањем задржао велики број танких свећица након затварања тржишта, чији
је медијан обима реда величине четрдесет пута мањи од обима у правој сесији.

Примењују се следећи кораци:

\begin{enumerate}
\item \textbf{Филтрирање према сесији.} Задржавају се само свећице унутар календарског
      прозора сесије за тај конкретан дан. Дани који нису права берзанска сесија
      одбацују се у целини, јер свећице од добављача повремено упадају на нерадне дане.
\item \textbf{Одбацивање артефаката.} Уклањају се свећице које истовремено имају нулти
      обим и нулту промену цене. Свећица са стварним обимом али без промене цене
      задржава се, јер представља легитиман миран тренутак.
\item \textbf{Уклањање дупликата} и сортирање по временској ознаци.
\item \textbf{Провера покривености.} Дан се задржава само ако број свећица достиже
      најмање 80\% очекиваног броја за тај конкретан дан и притом није мањи од
      апсолутног прага од двадесет опсервација.
\end{enumerate}

Два прага примењују се заједно јер ниједан сам није довољан: релативни праг од 80\% за
петнаестоминутну мрежу износи свега 20,8 свећица, па би дан са двадесет једном
свећицом прошао, док апсолутни праг искључује скраћене дане на тој мрежи, где број
свећица пада на четрнаест.

Праг се примењује \textbf{једном, на нивоу целе мреже}, а не засебно за сваки
естиматор. Примена по естиматору оставила би дизајне засноване на биповер варијацији
са више употребљивих дана него дизајне засноване на медијалној реализованој варијанси
на истој мрежи. Резултат би био неуравнотежен факторски план, што би уделе варијансе у
декомпозицији учинило неинтерпретабилним.

Поред тога, за сваки симбол спроводи се \textbf{тврда провера прилагођености
корпоративним акцијама}. Ако се утврди да серија цена није прилагођена поделама
акција, изградња панела се зауставља, јер би непрелагођена серија произвела лажне
скокове на датумима подела. Ако само преузимање података о корпоративним акцијама не
успе због мрежне грешке, провера се прескаче уз упозорење, али се процес не прекида.

\subsection{Конструкција панела}

За сваки симбол и сваку од тридесет шест комбинација дизајна мерења спроводи се
следећи ток:

\begin{enumerate}
\item преузимање и чишћење једноминутних свећица;
\item преузорковање на мрежу дефинисану димензијом А1 и рачунање логаритамских приноса;
\item рачунање реализованих мера и условне декомпозиције, према димензијама А2 и А3;
\item изградња HAR агрегата клизањем преко \textbf{доступних трговачких дана}, никада
      календарских, при чему се непотпуни прозори остављају као недостајући;
\item провера према замрзнутој шеми панела;
\item уписивање у колонски формат за даљу обраду.
\end{enumerate}

Прекноћни приноси ни у једном тренутку не улазе у реализоване мере: низ приноса сваког
дана гради се искључиво из цена тог истог дана.

Пре уписивања, над састављеним панелом спроводе се две тврде провере исправности:
идентитет семиваријанси из једначине~\eqref{eq:semivar-identity} и поштовање
апсолутног прага броја интрадневних опсервација. За разлику од бројаних дијагностика,
ове провере заустављају процес ако не прођу, јер њихов пад указује на грешку у рачуну,
а не на својство података.

\section{Дизајн експеримента}

\subsection{Фактори и нивои}

Експеримент укршта две експерименталне осе, чији су нивои дати у
табели~\ref{tab:nivoi}.

\begin{table}[htbp]
\centering
\caption{Експерименталне осе и њихови нивои}
\label{tab:nivoi}
\begin{tabularx}{\textwidth}{L{2.9cm} L{3.5cm} X c}
\toprule
Оса & Димензија или група & Нивои & Број \\
\midrule
Дизајн мерења
 & А1 --- фреквенција узорковања & 1m, 5m, 15m & 3 \\
 & А2 --- естиматор отпоран на скокове & BPV, MedRV, TrBPV & 3 \\
 & А3 --- тест за детекцију скокова & наивни, BNS, LM, LM-FDR & 4 \\
\addlinespace
Архитектура модела
 & економетријски & наивни, HAR, log-HAR, HAR-J, CHAR, SHAR, HARQ, HAR-IV, ARFIMA, MEM & 10 \\
 & засновани на стаблима & LightGBM, XGBoost & 2 \\
 & секвенцијални & PatchTST & 1 \\
 & темељни & ttm-rv, ttm-c, ttm-c+j, log-HAR-TTM & 4 \\
\bottomrule
\end{tabularx}
\end{table}

Комбиновањем трију димензија дизајна мерења добија се тридесет шест дизајна, а укупан
број модела износи седамнаест. Свака комбинација оцењује се на три хоризонта
($h = 1$, $h = 5$ и $h = 22$ дана) и кроз једанаест валидационих прозора.

\subsection{Свођење рачунског обима}

Пун декартов производ свих комбинација износио би
$36 \times 17 \times 3 \times 11 = 20\,196$ оцењивања, што је непотребно велико, јер
већина модела не зависи од свих трију димензија дизајна мерења. Модел који као улаз
користи само реализовану варијансу даје идентичне прогнозе под свих дванаест
комбинација естиматора и теста на истој мрежи, па би његово поновно оцењивање било
само трошење времена.

Сваки модел зато декларише \textbf{од којих димензија његови улази стварно зависе}, а
систем га оцењује једном по свакој различитој комбинацији само тих димензија, како је
приказано у табели~\ref{tab:zavisnost}.

\begin{table}[htbp]
\centering
\caption{Зависност модела од димензија дизајна мерења}
\label{tab:zavisnost}
\begin{tabularx}{\textwidth}{L{2.6cm} X c}
\toprule
Зависност & Модели & Оцењивања \\
\midrule
само А1 &
наивни, HAR, log-HAR, SHAR, HARQ, HAR-IV, ARFIMA, MEM, ttm-rv, log-HAR-TTM &
3 \\
\addlinespace
А1, А2 и А3 &
HAR-J, CHAR, LightGBM, XGBoost, PatchTST, ttm-c, ttm-c+j &
36 \\
\bottomrule
\end{tabularx}
\end{table}

Ова одлука није само рачунска. Понављање идентичних редова кроз дизајне од којих модел
не зависи унело би вештачки дуплиране опсервације у декомпозицију варијансе и исказало
би лажну прецизност.

\subsection{Фиксирани параметри}

Све вредности наведене у табели~\ref{tab:fiksirani} идентичне су у свим условима
експеримента и нису подешаване према резултатима.

\begin{table}[htbp]
\centering
\caption{Фиксирани параметри експеримента}
\label{tab:fiksirani}
\begin{tabularx}{\textwidth}{L{4.2cm} L{3.4cm} X}
\toprule
Параметар & Вредност & Разлог фиксирања \\
\midrule
Период узорка & 2010--2025 & расположивост једноминутне историје \\
Праг за TrBPV & 3,0 локалне девијације & увођење као четврте димензије дало би неуравнотежен план \\
Локални прозор LM теста & 4 сесије на свакој мрежи & да се са мрежом мења само фреквенција, а не и дужина прозора \\
Клизни прозор FDR корекције & 250 дана & приближно једна трговачка година; спречава гледање унапред \\
Минимална покривеност дана & 80\% очекиваног броја & филтрирање непотпуних сесија \\
Апсолутни праг опсервација & 20 & искључује петнаестоминутне скраћене дане \\
Хиперпараметри стабала & 200 стабала, дубина 4, 15 листова, стопа 0,05 & подешавање би мешало труд у претрази са ефектом архитектуре \\
Архитектура PatchTST & секвенца 60, сегмент 10, димензија 32, 4 главе, 2 слоја & намерно мала због величине узорка \\
Тежине хибридног модела & 0,5 и 0,5 & оптимизација би одразила претрагу, не архитектуру \\
Број семена & 5 & разликовање ефекта архитектуре од шума иницијализације \\
Контролна тачка темељног модела & тачно прикована ознака издања & различита издања имају различите скупове предтренирања \\
\bottomrule
\end{tabularx}
\end{table}

\subsection{Ток једног валидационог прозора}

Оцењивање сваког модела, за сваки дизајн и хоризонт, одвија се кроз једанаест
усидрених прозора описаних у одељку~\ref{subsec:validacija}. Један прозор пролази кроз
следеће фазе:

\begin{enumerate}
\item \textbf{Подела.} Из панела се издвајају тренинг период са фиксираним почетком,
      валидациони преклоп као последња година тренинга и тест период као наредна година.
\item \textbf{Оцењивање.} Модел се оцењује искључиво на тренинг делу, без валидационог
      преклопа.
\item \textbf{Оцена корекције.} На валидационом преклопу рачуна се непараметарска
      корекција повратне трансформације, као и коефицијенти регресије за накнадну
      рекалибрацију. Овај преклоп се никада не користи за подешавање хиперпараметара
      нити за евалуацију.
\item \textbf{Предвиђање.} Прогнозе се производе за тест период и враћају у ниво помоћу
      корекције оцењене у претходном кораку.
\item \textbf{Бележење.} Уписују се прогноза, вредност пре повратне трансформације,
      износ корекције, ознака прозора и све величине потребне за евалуацију.
\end{enumerate}

Прогнозе из свих једанаест прозора спајају се у јединствен скуп, над којим се потом
рачунају све метрике. Прозори се не усредњавају међусобно, јер би усредњавање по
прозору дало једнаку тежину прозорима различите дужине тест периода.

\subsection{Провера контаминације предтренирања}

Пре него што се произведе иједна прогноза темељног модела, спроводи се провера у којој
се објављени датум одсецања података за предтренирање пореди са тест периодом.
Резултат се уписује као трајан артефакт \textbf{у оба случаја}, и када преклапања нема
и када га има.

Када се контаминација утврди, преклапање се квантификује, а модел се задржава у
анализи уз изричиту ознаку да је контаминиран. Модел се никада не избацује тихо нити
се тихо третира као чист, јер би оба поступка учинила резултат непроверљивим.

\section{Мере}

Мере су унапред подељене у три групе, како би се избегло накнадно тражење статистички
значајних разлика међу великим бројем излазних величина.

\paragraph{Примарне мере.}
Служе за проверу хипотеза и наводе се у главним табелама:

\begin{itemize}
\item QLIKE функција губитка, као примарна мера тачности;
\item релативни QLIKE у односу на референтни модел, као приказна мера;
\item удели објашњене варијансе губитка по фактору, из декомпозиције варијансе;
\item припадност скупу поверења модела.
\end{itemize}

\paragraph{Секундарне мере.}
Служе за поређење са ранијом литературом и за проверу доследности:

\begin{itemize}
\item RMSE и MAE, уз обавезну варијанту MAE прилагођену условној медијани;
\item коефицијенти регресије реализоване вредности на прогнозу;
\item рекалибрисани QLIKE, који се наводи упоредо са сировим, никада уместо њега;
\item резултати тестова вредности при ризику.
\end{itemize}

\paragraph{Дијагностичке мере.}
Не служе за проверу хипотеза, него за контролу исправности:

\begin{itemize}
\item удео дана са детектованим скоком по дизајну;
\item индекс преклапања између варијанти теста за детекцију скокова;
\item удео скоковите компоненте у укупној варијанси;
\item број искључених дана по разлогу искључивања;
\item број опсервација по ћелији факторског плана.
\end{itemize}

\section{Правила искључивања}

Следећа правила утврђена су \textbf{пре увида у резултате}:

\begin{itemize}
\item дан који не задовољава критеријуме покривености искључује се за целу мрежу, а не
      само за поједини естиматор, како би факторски план остао уравнотежен;
\item дани у периоду загревања клизних прозора не улазе у статистику о скоковима, јер
      за њих локална оцена волатилности није дефинисана;
\item дани за које циљна променљива није дефинисана, дакле последњих $h$ дана сваке
      серије, искључују се из евалуације на том хоризонту;
\item симбол чија серија има мање од шездесет опсервација у тренинг делу прозора
      изоставља се из оцењивања модела који се оцењују засебно по симболу, уз бележење;
\item резултат се никада не искључује због тога што одступа од очекиваног. Дани са
      изразито великим губитком задржавају се у примарној анализи, а њихов утицај се
      проверава засебном варијантом декомпозиције из које је изузет горњи проценат дана
      по губитку.
\end{itemize}

\section{Статистичка анализа}

Провера хипотеза изведена је како следи.

\paragraph{Х1 --- доприноси ли информација о скоковима тачности.}
Пореде се модели који користе компоненте скокова са моделима који их не користе,
унутар истог дизајна и хоризонта. Значајност се утврђује упареним тестом предиктивне
тачности над серијама дневних губитака, уз корекцију за преклапајуће циљне вредности.

\paragraph{Х2 --- јача ли ефекат на краћем хоризонту.}
Иста поређења понављају се одвојено за сваки хоризонт, а закључак се изводи из
поређења величине ефекта, не из пуког броја значајних резултата.

\paragraph{Х3 --- колики је удео дизајна мерења у односу на избор модела.}
Ово је централна хипотеза. Дневни губици се разлажу декомпозицијом варијансе из
одељка~\ref{subsec:anova}, одвојено по хоризонту, са симболом као блокирајућим
фактором и уз троје структурних интеракција. Удели се исказују као дескриптивни удели
објашњене варијансе, а свака наведена значајност чита се уз двосмерно кластерисане
стандардне грешке.

\paragraph{Х4 --- мења ли се рангирање при преласку на економску меру.}
Рангирање по QLIKE пореди се са рангирањем по тестовима вредности при ризику. Тестови
се рачунају на обједињеној унији свих тест прозора, јер један прозор нема довољну
снагу.

\paragraph{Х5 --- опстаје ли побољшање уз насумично пермутоване индикаторе скокова.}
Индикатори се пермутују кроз време, чиме се задржава њихова учесталост а уништава веза
са стварним данима, након чега се поређење понавља.

Пошто се укупно спроводи велики број тестова, свака извештавана породица поређења
подвргава се контроли стопе лажних открића. Породица се наводи експлицитно и бележи уз
резултат, јер контрола нема смисла без изјаве над чим се спроводи. Уз сваку наведену
вредност значајности наводи се и величина ефекта, а закључак се никада не изводи само
из вредности значајности.

\subsection{Робусност}

Декомпозиција варијансе спроводи се у четири варијанте наведене у
табели~\ref{tab:robusnost}, једна поред друге. Стабилност удела кроз све четири
варијанте представља основ за тврдњу о робусности, док се нестабилност наводи као
налаз, а не прикрива.

\section{Рачунско окружење и поновљивост}

Цео поступак имплементиран је у програмском језику Python, уз управљање зависностима
које фиксира тачне верзије свих коришћених библиотека.

Две особине имплементације од методолошког су значаја. Прва је
\textbf{регистарски образац}: естиматори, тестови за детекцију скокова и модели
разрешавају се искључиво преко назива, из централног регистра, па се додавање нове
варијанте своди на једну нову датотеку, без измене иједног постојећег фајла. Тиме је
структурно онемогућено да било која варијанта добије повлашћен третман у коду, што је
предуслов за поштено поређење. Друга је \textbf{детерминизам}: једини стохастички
модел оцењује се на пет унапред задатих семена, а примарна извештавана величина
усредњава губитке по семену, не прогнозе, док ансамбл усредњених прогноза постоји као
засебна, секундарна величина.

Аутоматизовани тестови покривају критичне делове: генерисање валидационих прозора,
исправност повратне трансформације, поштовање регистарског обрасца и одсуство преласка
између симбола при грађењу циљне променљиве и секвенци.

Подаци се не чувају уз изворни код, јер су у потпуности поново изводиви из извора на
основу забележених манифеста порекла.

\section{Ограничења методологије}

Следећа ограничења проистичу из самог дизајна истраживања и наводе се унапред.

\begin{itemize}
\item \textbf{Зависност између дизајна.} Свих тридесет шест дизајна изведено је из
      истих сирових једноминутних свећица, па су вредности реализоване варијансе за
      исти дан високо корелисане. Претпоставка независности класичног теста стога не
      важи, због чега се удели варијансе исказују као дескриптивни, а свака значајност
      чита уз кластерисане стандардне грешке.
\item \textbf{Једно тржиште.} Сви симболи деле исти трговачки календар, па рад не
      проверава да ли закључци важе на тржишту са другачијим радним временом.
\item \textbf{Искључен прекноћни принос.} Реализована волатилност обухвата само
      интрадневно кретање, па се и економска евалуација односи искључиво на интрадневни
      ризик.
\item \textbf{Изостављени модели.} Скуп модела не обухвата све архитектуре из
      литературе; избор је ограничен на оне које су или стандард у области или
      средиште актуелне расправе.
\item \textbf{Фиксирани хиперпараметри.} Модели засновани на стаблима и секвенцијални
      модел нису подешавани, па њихови резултати представљају доњу границу могућег
      учинка тих архитектура. Та одлука је намерна, јер би подешавање помешало уложен
      труд у претрагу са ефектом саме архитектуре.
\end{itemize}
